\documentclass{article}
\usepackage[UTF8]{ctex}
\usepackage{amsmath,amssymb,amsfonts}
\usepackage{enumitem}
\usepackage{geometry}
\geometry{a4paper, margin=1in}
% === 错题本排版优化（不参与题目标识） ===
\setlength{\parindent}{0pt}
\setlength{\parskip}{0.35em plus 0.10em minus 0.08em}
\setlength{\emergencystretch}{3em}
\setlist{itemsep=0.25em, topsep=0.40em, parsep=0pt}
\allowdisplaybreaks[2]
% === 错题本排版优化结束 ===
\begin{document}
\section*{东北大学2026年高等代数真题}
\section*{67．东北大学2026年高等代数试卷}

\subsection{题目1}
1．（15 分）计算 \(n\) 阶行列式

\[
D_{n}=\left|\begin{array}{ccccc}
2 & & & & 2 \\
-1 & 2 & & & 2 \\
& \ddots & \ddots & & \vdots \\
& & -1 & 2 & 2 \\
& & & -1 & 2
\end{array}\right| \text { (用两种方法). }
\]

\subsection{题目2}
2．（15 分）设线性方程组

\[
\left\{\begin{array}{l}
x_{1}+\lambda x_{2}+\mu x_{3}+x_{4}=0 \\
2 x_{1}+x_{2}+x_{3}+2 x_{4}=0 \\
3 x_{1}+(2+\lambda) x_{2}+(4+\mu) x_{3}+4 x_{4}=1
\end{array}\right.
\]

其中 \((1,-1,1,-1)^{\prime}\) 为方程的解．

\begin{itemize}
  \item[（1）] 求该方程组的通解。
  \item[（2）] 求符合 \(x_{2}=x_{3}\) 的所有解．
\end{itemize}

\subsection{题目3}
3．（15 分）解答如下问题：

\begin{itemize}
  \item[（1）] 求 \(x^{4}+x^{2}+1=0\) 的所有复根．
  \item[（2）] 设 \(f_{1}(x), f_{2}(x)\) 是次数不超过3的首一互异多项式，且
\[
x^{4}+x^{2}+1 \mid f_{1}\left(x^{3}\right)+x^{4} f_{2}\left(x^{3}\right)
\]
求 \(\left(f_{1}(x), f_{2}(x)\right)\) ．
\end{itemize}

\subsection{题目4}
4．（15 分）解答如下问题：

\begin{itemize}
  \item[（1）] 设 \(\varepsilon_{1}, \varepsilon_{2}, \varepsilon_{3}\) 是线性空间 \(V\) 的一组基，线性变换 \(\mathcal{A}\) 在 \(\varepsilon_{1}, \varepsilon_{2}, \varepsilon_{3}\) 下的矩阵为
\[
A=\left(\begin{array}{ccc}
0 & 3 & -1 \\
1 & -2 & 2 \\
4 & 1 & -1
\end{array}\right)
\]
求 \(\mathcal{A}\left(2 \varepsilon_{1}-\varepsilon_{2}+5 \varepsilon_{3}\right)\) 在 \(\varepsilon_{1}, \varepsilon_{2}, \varepsilon_{3}\) 下的坐标．
  \item[（2）] 设 \(\alpha\) 为 \(n\) 维线性空间 \(V\) 中的非零向量，\(V\) 上的线性变换 \(\mathcal{A}\) 满足 \(\mathcal{A}^{n-1} \alpha \neq 0, \mathcal{A}^{n} \alpha=0\) ，证明：\(a, \mathcal{A} \alpha, \cdots, \mathcal{A}^{n-1} a\) 构成 \(V\) 的一组基．
\end{itemize}

\subsection{题目5}
5．（15 分）设 \(U=L\left(\alpha_{1}, \alpha_{2}\right)\) 为欧氏空间 \(\mathbb{R}^{4}\) 的子空间，其中 \(\alpha_{1}=(1,1,2,1)^{\prime}, \alpha_{2}=(1,0,0,-2)^{\prime}\) ，取标准内积，即 \((\alpha, \beta)=\alpha^{\prime} \beta, \alpha, \beta \in \mathbb{R}^{4}\) ，求 \(\operatorname{dim} U^{\perp}\) 和 \(U^{\perp}\) 的一组标准正交基。

\subsection{题目6}
6．（15 分）设 \(\alpha_{1}, \alpha_{2}, \alpha_{3}, \alpha_{4}\) 为线性空间 \(V\) 的一组基，线性变换 \(\mathcal{A}\) 在 \(\alpha_{1}, \alpha_{2}, \alpha_{3}, \alpha_{4}\) 下的矩阵为

\[
A=\left(\begin{array}{cccc}
1 & 0 & 2 & 1 \\
-1 & 2 & 1 & 3 \\
1 & 2 & 5 & 5 \\
2 & -2 & 1 & -2
\end{array}\right) .
\]

求 \(\operatorname{Ker} \mathcal{A}, \operatorname{Im} \mathcal{A}\) ．

\subsection{题目7}
7．（15 分）设 \(A=\left(\begin{array}{lll}1 & 0 & 0 \\ 1 & 0 & 1 \\ 0 & 1 & 0\end{array}\right)\)

\begin{itemize}
  \item[（1）] 证明：\(A^{n}=A^{n-2}+A^{2}-E(n \geq 3)\) ．
  \item[（2）] 证明：\(A^{100}-E\) 不是零矩阵．
\end{itemize}

\subsection{题目8}
8．（15 分）设 \(A=\left(\begin{array}{ll}a & b \\ c & d\end{array}\right) \in \mathbb{R}^{2 \times 2},|A|=1\) ．证明：

\begin{itemize}
  \item[（1）] 当 \(|a+d|>2\) 时，\(A\) 与 \(\left(\begin{array}{cc}\lambda & 0 \\ 0 & \frac{1}{\lambda}\end{array}\right)\) 相似，其中 \(\lambda \neq \pm 1,0\) 且 \(\lambda \in \mathbb{R}\) ．
  \item[（2）] 当 \(|a+d|=2\) 时，则 \(A= \pm\left(\begin{array}{ll}1 & 0 \\ 0 & 1\end{array}\right)\) 或相似于 \(\left(\begin{array}{ll}1 & 1 \\ 0 & 1\end{array}\right)\) 或 \(\left(\begin{array}{cc}-1 & 1 \\ 0 & -1\end{array}\right)\) ．
  \item[（3）] 当 \(|a+d|<2\) 时，证明：存在可逆矩阵 \(P\) ，使得 \(P^{-1} A P=\left(\begin{array}{cc}\cos \theta & \sin \theta \\ -\sin \theta & \cos \theta\end{array}\right)\) ，并求 \(P\) ．
\end{itemize}

\subsection{题目9}
9．（15 分）设 \(\alpha_{1}, \alpha_{2}, \cdots, \alpha_{n}\) 为线性空间 \(V\) 的一组基，线性变换 \(\mathcal{A}\) 在 \(\alpha_{1}, \alpha_{2}, \cdots, \alpha_{n}\) 下的矩阵为

\[
\left(\begin{array}{cccc}
a & 1 & & \\
& a & & \\
& & & 1 \\
& & & a
\end{array}\right)
\]

（1）设 \(\alpha_{n}\) 为 \(\mathcal{A}-\) 子空间 \(W\) 的一个向量，证明：\(W=V\) ．\\
（2）证明：\(\alpha_{1}\) 属于所有非零的 \(\mathcal{A}-\) 子空间。\\
（3）证明：\(V\) 不能表示为两个非平凡的 \(\mathcal{A}-\) 子空间的直和．

\subsection{题目10}
10．（15 分）设 \(A, B\) 都为 \(n\) 阶方阵。

\begin{itemize}
  \item[（1）] 证明：\(A B, B A\) 有相同的特征值．
  \item[（2）] 证明：不存在正交矩阵 \(A, B\) ，使得 \(A^{2}=A B+B^{2}\) ．
\end{itemize}
\end{document}
